\section{Forests whose terminal vertices are all far from one port}
\label{sec:far-terminal-return-forests}

The distance-layer encoding also controls several paths whose endpoints all lie in the far terminal tail.
Only one of those paths can traverse any ordinary two-vertex layer.
This avoids the linear factor needed for two cycles with independently placed layer intervals.

\begin{theorem}[A fixed number of far terminal vertices]
\label{thm:far-terminal-forests}
Fix an even integer $q\ge2$ and $0<\varepsilon\le1/2$.
There is a constant $R_{q,\varepsilon}$ with the following property.
Let $T$ be a connected loopless multigraph of rank $h$ with distinct degree-two ports
\[
 c,a_1,\ldots,a_q.
\]
Every other vertex has degree three.
Suppose
\[
 P_{ca_i}(T)\ge(1/2+\varepsilon)2^h
 \qquad(1\le i\le q).
\]
Then the number of forests whose odd-degree vertex set is exactly $\{a_1,\ldots,a_q\}$ is at most
\[
 R_{q,\varepsilon}2^{h/2}.
\]
Each such forest consists of $q/2$ vertex-disjoint paths pairing the leaf ports.
\end{theorem}

\begin{proof}
Put $C=C_\varepsilon$ for the additive path-stability constant.
Suppress all ports except $c,a_i$ when considering that terminal pair.
No loop can arise.
A loop would represent a cyclic part avoided by every terminal path and would force normalized density at most one half.
The resulting normalized two-pole has rank $h$ and the same terminal-path count.
Additive stability gives terminal distance at least $h-C$ there, and hence also in $T$.
Thus
\begin{equation}\label{eq:far-port-distances}
 \operatorname{dist}_T(c,a_i)\ge h-C
 \qquad(1\le i\le q).
\end{equation}

Assume first that $T$ is bridgeless.
Suppress $a_2,\ldots,a_q$ and call the resulting normalized two-pole $N$.
It is bridgeless and has order $2h$.
Choose a shortest $c$--$a_1$ path $A$ in $N$.
The residual-rank part of additive stability gives
\[
 H:=\beta(N-V(A))\le C.
\]

Each edge of $N$ corresponds to a path of $T$, with at most $q-1$ suppressed vertices in total.
Consider an edge whose path contains a suppressed port, and one endpoint $v$ of that edge.
A shortest $c$--$v$ path in $N$ lifts to a walk of length at most
$\operatorname{dist}_N(c,v)+q-1$ in $T$.
Following the subdivided edge from $v$ reaches that port in at most $q$ further steps.
Equation~\eqref{eq:far-port-distances} therefore implies
\[
 \operatorname{dist}_N(c,v)\ge h-C-2q.
\]
The slightly weaker constant avoids distinctions about overlap between these two walks.

Ranks bounded in terms of $q,C$ satisfy the theorem by the trivial bound $2^h$.
For the remaining ranks, put
\[
 d=\lfloor h-C-2q-2\rfloor\ge1,
 \qquad \Delta=h-d\le C+2q+3.
\]
Use the distance layers of $N$ from $c$.
All suppressed subdivision paths have both endpoints in layers at least $d$.
Also $a_1$ lies in that tail, since its distance in $N$ is at least $h-C$.

Mark $c$, every vertex in a layer at least $d$, and every internal layer with at least three vertices.
As in the two-pole layer proof, at most $6\Delta+2$ vertices of $N$ are marked.
Every remaining layer has exactly two vertices, one on $A$.
Mark whole layers containing endpoints of parallel pairs outside $A$.
This adds at most $4H$ vertices, since their digons are independent in $N-V(A)$.
Restore the $q-1$ suppressed ports and mark them as well.
Every restored edge segment lies wholly in the marked subgraph.
Outside the marked set, the graphs $T$ and $N$ agree.
The marked set has size bounded in terms of $q,\varepsilon$ alone.

Let $F$ be one forest counted in the theorem.
At each leaf port its degree is one.
At every other vertex its degree is zero or two.
Therefore it consists of $q/2$ vertex-disjoint paths, with all endpoints in the marked terminal tail.

If $F$ lies wholly in the tail, all its edges are determined by marked incidences.
Otherwise let $i<d$ be its minimum distance layer among the vertices of $N$.
If this layer is ordinary, mark its two vertices.
This adds at most two vertices to the marked set $B$.
Fix all forest edges incident with $B$, with at most $2^{3|B|}$ choices.

Choose a path component of $F$ that attains layer $i$.
Both its endpoints lie in the terminal tail.
It therefore traverses every ordinary layer strictly between $i$ and $d$ with two strands.
Both vertices in each such layer belong to this path, with one backward and one forward edge.
No other component of $F$ can use either vertex.
Consequently, the entire forest has exactly the same matching description on ordinary layers as one returning path.

Every cut between consecutive occupied ordinary layers has at most two matching choices.
Two-choice cuts are nonconsecutive.
The extra parallel-edge markings justify the same assertion in a multigraph.
The matching choices contribute at most $2^{(d-i)/2}$.
The remaining edges are determined by the marked-incidence data.
Hence
\[
 \#F\le O_{q,\varepsilon}(1)
 \left(1+\sum_{i=0}^{d-1}2^{(d-i)/2}\right)
 =O_{q,\varepsilon}(2^{d/2})
 =O_{q,\varepsilon}(2^{h/2}).
\]

For completeness, consider a bridge in $T$.
A shore with $n$ vertices and $p$ original ports has rank $(n-p+1)/2$, a positive integer.
Thus both shores contain cycles.
Every high terminal pair must cross the bridge, since otherwise it avoids a cyclic shore and has normalized density at most one half.
It follows that every bridge separates $c$ from all $a_i$.
No original port can be a bridge endpoint, since its terminal paths would immediately leave a cyclic shore.

Delete one bridge and retain its leaf-port shore.
Its new bridge attachment serves as the center port.
All normalized center-to-leaf path densities remain at least $1/2+\varepsilon$, by factorization across the bridge.
Every counted forest stays on this shore.
Indeed, each path has both endpoints there and cannot cross a bridge twice.
The shore has smaller rank.
Induction on rank proves the result with the same constant.
\end{proof}

\begin{corollary}[The four far ports of a five-pole]
\label{cor:five-pole-four-far-forests}
Let $T$ be a normalized five-pole of rank $h$, with center port $c$ and four other ports.
If every center-to-other-port path count is at least $(1/2+\varepsilon)2^h$, then the total forest count on those four other ports is
\[
 O_\varepsilon(2^{h/2}).
\]
In particular, this applies to a full high star in a five-pole.
\end{corollary}

This lemma is an input for a possible five-edge-cut state analysis.
It does not by itself prove a half-exponent bound for all five-edge cuts.
